% !PRINTED 80
\weberpage{80}

From this, however, one obtains:

\weq{20}{
\begin{aligned}
v_1' &= \tfrac{1}{2}(v_1+v_2+v_3+v_4), & a_1 &= \tfrac{1}{2}(\phantom{-}v_2+v_3+v_4),\\
v_2' &= \tfrac{1}{2}(v_1+v_2-v_3-v_4), & a_2 &= \tfrac{1}{2}(\phantom{-}v_2-v_3-v_4),\\
v_3' &= \tfrac{1}{2}(v_1-v_2+v_3-v_4), & a_3 &= \tfrac{1}{2}(-v_2+v_3-v_4),\\
v_4' &= \tfrac{1}{2}(v_1-v_2-v_3+v_4), & a_4 &= \tfrac{1}{2}(-v_2-v_3+v_4),
\end{aligned}
}

and formula (18) then gives:
\[
4\,A_1\,\vartheta_{00}(v_1) = \sum \vartheta_{g_1,g_2}(v_1')\vartheta_{g_1,g_2}(v_2')\vartheta_{g_1,g_2}(v_3')\vartheta_{g_1,g_2}(v_4'),
\]
in which now either the $v_i$ or the $v_i'$ may be regarded as independent variables.

If one regards the $v_i$ as independent variables, then in this formula $A_1$ is independent of $v_1$, but still depends on $v_2$, $v_3$, $v_4$. If we therefore set $4\,A_1 = c\,\vartheta_{00}(v_2)\vartheta_{00}(v_3)\vartheta_{00}(v_4)$, there results
\[
c\,\vartheta_{00}(v_1)\vartheta_{00}(v_2)\vartheta_{00}(v_3)\vartheta_{00}(v_4)
= \sum \vartheta_{g_1,g_2}(v_1')\vartheta_{g_1,g_2}(v_2')\vartheta_{g_1,g_2}(v_3')\vartheta_{g_1,g_2}(v_4'),
\]
and in this $c$ is in any case independent of $v_1$. But since the right-hand side of this formula remains unchanged under arbitrary permutations of $v_1, v_2, v_3, v_4$, $c$ is independent of all the $v_i$, and one obtains its value by setting all $v_i = 0$:
\[
c\,\vartheta_{00}^4 = \vartheta_{00}^4 + \vartheta_{01}^4 + \vartheta_{10}^4,
\]
from which, by \S\,21, (14), it follows that $c=2$.

We thus have the formula:
\[
2\,\vartheta_{00}(v_1)\vartheta_{00}(v_2)\vartheta_{00}(v_3)\vartheta_{00}(v_4)
= \sum\limits^{g_1,g_2} \vartheta_{g_1,g_2}(v_1')\vartheta_{g_1,g_2}(v_2')\vartheta_{g_1,g_2}(v_3')\vartheta_{g_1,g_2}(v_4').
\]

If in it each of the variables $v_1, v_2, v_3, v_4$ is increased by a half period $-\tfrac{1}{2}(g_1'\omega_1 - g_2')$, then $v_1'$ is increased by a whole period $-(g_1'\omega_1 - g_2')$, while $v_2', v_3', v_4'$ remain unchanged, and accordingly one obtains, with the help of \S\,21, (2) and (3), a system of four formulae:
% [sic] German source prints \omega_1 in both shifts; in the normalized theta setting the period is \omega.

\weq{21}{
\begin{aligned}
&2\,\vartheta_{g_1',g_2'}(v_1)\vartheta_{g_1',g_2'}(v_2)\vartheta_{g_1',g_2'}(v_3)\vartheta_{g_1',g_2'}(v_4)\\
&=\sum\limits^{g_1,g_2} (-1)^{g_1g_2'+g_2g_1'+g_1'g_2'}\,\vartheta_{g_1,g_2}(v_1')\vartheta_{g_1,g_2}(v_2')\vartheta_{g_1,g_2}(v_3')\vartheta_{g_1,g_\nu}(v_4'). % [sic] German source prints g_\nu; mathematically expected g_2
\end{aligned}
}

These are the \gesperrt{Jacobi} formulae.
If one sets, for example,
\[
\begin{aligned}
v_1 &= 0, & v_2 &= v+w, & v_3 &= u+v, & v_4 &= u+w,\\
v_1' &= u+v+w, & v_2' &= -u, & v_3' &= -w, & v_4' &= -v,
\end{aligned}
\]
%CONT>>
